% Abhakseite „Das kann ich“ – Zone nur im Gesamt (\mitzone)
%% TODO Ich-kann-Sätze: Platzhalter wie in den Aufgabendateien
\begin{abhakseite}
\ifdefined\mitzone
\abhakgruppe{Kennst du schon}
\abhak{1}{Eine Gleichung nach x auflösen und die Variablen tauschen, bei e-Funktionen mit ln (Exponentialgleichung)}
\abhak{2}{Wertebereich einer Funktion aus Monotonie und Grenzverhalten angeben, strenge Monotonie am Graphen oder über die Ableitung erkennen}
\abhak{3}{Punkte an einer Geraden spiegeln, Bildpunkte benennen, Spiegelachse als Symmetrieachse einer Figur}
\abhak{4}{Integral einer Differenz als Fläche zwischen zwei Graphen deuten, Trapezformel und gleichschenklig-rechtwinkliges Dreieck}
\abhak{5}{Steigung einer Geraden und Steigung der gespiegelten Geraden (Kehrwert; Steigung minus eins senkrecht zu y = x)}
\abhak{6}{Wertebereich einer Funktion aus Monotonie und Grenzverhalten angeben, strenge Monotonie am Graphen oder über die Ableitung erkennen – Fehler finden}
\abhak{7}{Wertebereich einer Funktion aus Monotonie und Grenzverhalten angeben, strenge Monotonie am Graphen oder über die Ableitung erkennen – selbst rechnen}
\fi
\abhakgruppe{Einheit 1 · Umkehrbarkeit, Bereiche und Term}
\abhak{8}{Umkehrbarkeit, Bereiche und Term}
\abhak{9}{Radius der Flüssigkeitsoberfläche in Abhängigkeit von der Füllhöhe über die Umkehrung der Profilfunktion nachweisen}
\abhak{10}{Umkehrbarkeit, Bereiche und Term – Fehler finden, Begründen, Anwendung, Darstellungswechsel}
\abhakgruppe{Einheit 2 · Spiegelung an der Winkelhalbierenden}
\abhak{11}{Spiegelung an der Winkelhalbierenden}
\abhak{12}{Spiegelung an der Winkelhalbierenden – Fehler finden, Begründen, Anwendung, Darstellungswechsel}
\end{abhakseite}
